% GENERATED FILE. Edit canonical lessons, then run npm run generate:books.
% canonical-source-hash: fnv1a64:c1d66bf8751ac1c3
% canonical-lessons: JA-C116-masu, JA-C116-funa, JA-C116-hamachi, JA-C116-sanma, JA-C116-asari

\chapter{Trout, Crucian carp, Young yellowtail, Pacific saury, Short-neck clam}
\label{ch:ja-trout-crucian-carp-young-yellowtail-pacific-saury-short-neck-clam}
\hlchaptermodality{116}

\begin{chapteropening}
\textbf{By the end of this chapter:} \emph{I can say a trout, a crucian carp, a young yellowtail, a Pacific saury and a short-neck clam.}

Complete the last lesson of chapter 116: i can say a trout, a crucian carp, a young yellowtail, a Pacific saury and a short-neck clam.
\end{chapteropening}

\section[\texorpdfstring{\ja{ます} (masu)}{ます (masu)}]{\ja{ます} (masu) — a trout}

\label{lesson:JA-C116-masu}



\begin{quote}
Before the new one: say the Japanese for a sea cucumber, then the Japanese for a sweetfish.
\end{quote}

\subsection*{You'll want to know: \ja{ます}}
\textbf{\ja{ます}} — \emph{masu} — \textquotedblleft{}a trout\textquotedblright{}.

\begin{grammarlens}[title={What you've built}]
One more word for this chapter.
\end{grammarlens}

\subsection*{Your turn}
\begin{itemize}
\raggedright
  \item \emph{Say it:} \emph{masu}
  \item \emph{Say it:} \emph{masu}, once more
  \item \emph{Say it:} say \emph{masu}
  \item \emph{Recall:} say the Japanese for a sea cucumber, then the Japanese for a sweetfish, then say \emph{masu} again
\end{itemize}

\begin{tcolorbox}[breakable,colback=teal!4,colframe=teal!35!black,title={Before you move on}]
What does \ja{ます} mean? (\textquotedblleft{}A trout\textquotedblright{}.) Say it once more.
\end{tcolorbox}
\section[\texorpdfstring{\ja{ふな} (funa)}{ふな (funa)}]{\ja{ふな} (funa) — a crucian carp}

\label{lesson:JA-C116-funa}



\begin{quote}
Before the new one: say the Japanese for a sweetfish, then the Japanese for a trout.
\end{quote}

\subsection*{You'll want to know: \ja{ふな}}
\textbf{\ja{ふな}} — \emph{funa} — \textquotedblleft{}a crucian carp\textquotedblright{}.

\begin{grammarlens}[title={What you've built}]
One more word for this chapter.
\end{grammarlens}

\subsection*{Your turn}
\begin{itemize}
\raggedright
  \item \emph{Say it:} \emph{funa}
  \item \emph{Say it:} \emph{funa}, once more
  \item \emph{Say it:} say \emph{funa}
  \item \emph{Recall:} say the Japanese for a sweetfish, then the Japanese for a trout, then say \emph{funa} again
\end{itemize}

\begin{tcolorbox}[breakable,colback=teal!4,colframe=teal!35!black,title={Before you move on}]
What does \ja{ふな} mean? (\textquotedblleft{}A crucian carp\textquotedblright{}.) Say it once more.
\end{tcolorbox}
\section[\texorpdfstring{\ja{はまち} (hamachi)}{はまち (hamachi)}]{\ja{はまち} (hamachi) — a young yellowtail}

\label{lesson:JA-C116-hamachi}



\begin{quote}
Before the new one: say the Japanese for a trout, then the Japanese for a crucian carp.
\end{quote}

\subsection*{You'll want to know: \ja{はまち}}
\textbf{\ja{はまち}} — \emph{hamachi} — \textquotedblleft{}a young yellowtail\textquotedblright{}.

\begin{grammarlens}[title={What you've built}]
One more word for this chapter.
\end{grammarlens}

\subsection*{Your turn}
\begin{itemize}
\raggedright
  \item \emph{Say it:} \emph{hamachi}
  \item \emph{Say it:} \emph{hamachi}, once more
  \item \emph{Say it:} say \emph{hamachi}
  \item \emph{Recall:} say the Japanese for a trout, then the Japanese for a crucian carp, then say \emph{hamachi} again
\end{itemize}

\begin{tcolorbox}[breakable,colback=teal!4,colframe=teal!35!black,title={Before you move on}]
What does \ja{はまち} mean? (\textquotedblleft{}A young yellowtail\textquotedblright{}.) Say it once more.
\end{tcolorbox}
\section[\texorpdfstring{\ja{さんま} (sanma)}{さんま (sanma)}]{\ja{さんま} (sanma) — a Pacific saury}

\label{lesson:JA-C116-sanma}



\begin{quote}
Before the new one: say the Japanese for a crucian carp, then the Japanese for a young yellowtail.
\end{quote}

\subsection*{You'll want to know: \ja{さんま}}
\textbf{\ja{さんま}} — \emph{sanma} — \textquotedblleft{}a Pacific saury\textquotedblright{}.

\begin{grammarlens}[title={What you've built}]
One more word for this chapter.
\end{grammarlens}

\subsection*{Your turn}
\begin{itemize}
\raggedright
  \item \emph{Say it:} \emph{sanma}
  \item \emph{Say it:} \emph{sanma}, once more
  \item \emph{Say it:} say \emph{sanma}
  \item \emph{Recall:} say the Japanese for a crucian carp, then the Japanese for a young yellowtail, then say \emph{sanma} again
\end{itemize}

\begin{tcolorbox}[breakable,colback=teal!4,colframe=teal!35!black,title={Before you move on}]
What does \ja{さんま} mean? (\textquotedblleft{}A Pacific saury\textquotedblright{}.) Say it once more.
\end{tcolorbox}
\section[\texorpdfstring{\ja{あさり} (asari)}{あさり (asari)}]{\ja{あさり} (asari) — a short-neck clam}

\label{lesson:JA-C116-asari}



\begin{quote}
Before the new one: say the Japanese for a young yellowtail, then the Japanese for a Pacific saury.
\end{quote}

\subsection*{You'll want to know: \ja{あさり}}
\textbf{\ja{あさり}} — \emph{asari} — \textquotedblleft{}a short-neck clam\textquotedblright{}.

\begin{grammarlens}[title={What you've built}]
One more word for this chapter.
\end{grammarlens}

\subsection*{Your turn}
\begin{itemize}
\raggedright
  \item \emph{Say it:} \emph{asari}
  \item \emph{Say it:} \emph{asari}, once more
  \item \emph{Say it:} say \emph{asari}
  \item \emph{Recall:} say the Japanese for a young yellowtail, then the Japanese for a Pacific saury, then say \emph{asari} again
\end{itemize}

\begin{tcolorbox}[breakable,colback=teal!4,colframe=teal!35!black,title={Before you move on}]
What does \ja{あさり} mean? (\textquotedblleft{}A short-neck clam\textquotedblright{}.) Say it once more.
\end{tcolorbox}
